\section{Part 0: The Postulates}

\subsection{Basic}

The differential substrate is modeled algebraically through a commutative ring carrying $n$ commuting derivations, which represent partial differentiation in an $n$-dimensional affine coordinate patch. This approach avoids analytic assumptions and allows every differential identity (Christoffel symbols, Riemann tensor, Weyl curvature) to be derived as purely algebraic consequences of the Leibniz rule and commutativity of mixed partials.

\begin{definition}[\lean{Basic.ScaleAlgebra}]
A \emph{scale algebra} of dimension $n$ on a commutative ring $A$ is a pair of data $(\mathbf{d}, A)$ where:
\begin{enumerate}
\item $\mathbf{d} = (\partial_0, \ldots, \partial_{n-1})$ is an $n$-tuple of endomorphisms $\partial_i : A \to A$;
\item each $\partial_i$ satisfies the Leibniz rule: for all $a, b \in A$,
  $$\partial_i(a \cdot b) = \partial_i(a) \cdot b + a \cdot \partial_i(b);$$
\item the mixed partials commute: for all $i, j$ and $a \in A$,
  $$\partial_i(\partial_j(a)) = \partial_j(\partial_i(a)).$$
\end{enumerate}
\end{definition}

\begin{theorem}[\lean{Basic.ScaleAlgebra.d\_zero}]
For any scale algebra and any $i \in \{0,\ldots,n-1\}$, the zero element is annihilated: $\partial_i(0) = 0$.
\end{theorem}
\begin{proof}
By the Leibniz rule applied to $0 + 0 = 0$, we have $\partial_i(0 + 0) = \partial_i(0) + \partial_i(0)$. Since $0 + 0 = 0$, this gives $\partial_i(0) = 2 \cdot \partial_i(0)$, which implies $\partial_i(0) = 0$ by cancellation.
\end{proof}

\begin{theorem}[\lean{Basic.ScaleAlgebra.d\_one}]
For any scale algebra and any $i \in \{0,\ldots,n-1\}$, the unity is annihilated: $\partial_i(1) = 0$.
\end{theorem}
\begin{proof}
By the Leibniz rule applied to $1 \cdot 1 = 1$, we have
$$\partial_i(1) = \partial_i(1) \cdot 1 + 1 \cdot \partial_i(1) = 2 \cdot \partial_i(1).$$
Thus $\partial_i(1) = 0$ by cancellation.
\end{proof}

\begin{definition}[\lean{Basic.ScaleAlgebra.dHom}]
Each derivation $\partial_i$ extends to an additive group homomorphism $\mathbf{d}_i : A \to A$, packaged as the morphism $(A, +) \to (A, +)$ that preserves zero and addition.
\end{definition}

\begin{theorem}[\lean{Basic.ScaleAlgebra.d\_neg}]
For any scale algebra and any $i$ and $a \in A$,
$$\partial_i(-a) = -\partial_i(a).$$
\end{theorem}
\begin{proof}
This is immediate from the fact that $\partial_i$ is an additive group homomorphism.
\end{proof}

\begin{theorem}[\lean{Basic.ScaleAlgebra.d\_sub}]
For any scale algebra and any $i$ and $a, b \in A$,
$$\partial_i(a - b) = \partial_i(a) - \partial_i(b).$$
\end{theorem}
\begin{proof}
This is immediate from the fact that $\partial_i$ is an additive group homomorphism.
\end{proof}

\begin{theorem}[\lean{Basic.ScaleAlgebra.d\_sum}]
For any scale algebra, any finite index set $\iota$, any function $f : \iota \to A$, and any $i \in \{0,\ldots,n-1\}$,
$$\partial_i\left(\sum_{x \in s} f(x)\right) = \sum_{x \in s} \partial_i(f(x)),$$
where $s$ is the underlying finite set of indices.
\end{theorem}
\begin{proof}
This follows from the fact that $\partial_i$ is an additive group homomorphism and preserves finite sums.
\end{proof}

\begin{theorem}[\lean{Basic.ScaleAlgebra.d\_natCast}]
For any scale algebra, any $i \in \{0,\ldots,n-1\}$, and any natural number $k$, the derivation annihilates the image of the natural number under the ring homomorphism: $\partial_i(k) = 0$.
\end{theorem}
\begin{proof}
By induction on $k$: the base case $\partial_i(0) = 0$ is \lean{Basic.ScaleAlgebra.d\_zero}. For the successor, $\partial_i(k+1) = \partial_i(k) + \partial_i(1) = 0 + 0 = 0$ by the Leibniz rule and inductive hypothesis.
\end{proof}

\begin{theorem}[\lean{Basic.ScaleAlgebra.d\_nsmul}]
For any scale algebra, any $i$, any natural number $k$, and any $a \in A$,
$$\partial_i(k \cdot a) = k \cdot \partial_i(a).$$
\end{theorem}
\begin{proof}
By the Leibniz rule, $\partial_i(k \cdot a) = \partial_i(k) \cdot a + k \cdot \partial_i(a) = 0 + k \cdot \partial_i(a)$ since $\partial_i(k) = 0$.
\end{proof}

The Kronecker delta is the bare Euclidean metric tensor, valued in $A$. It is pure bookkeeping and carries no physics.

\begin{definition}[\lean{Basic.kron}]
The Kronecker delta with indices $i, j \in \{0,\ldots,n-1\}$ is defined by
$$\delta_{ij} := \begin{cases} 1 & \text{if } i = j \\ 0 & \text{if } i \ne j. \end{cases}$$
\end{definition}

\begin{theorem}[\lean{Basic.kron\_self}]
For any $i \in \{0,\ldots,n-1\}$, $\delta_{ii} = 1$.
\end{theorem}
\begin{proof}
Immediate from the definition.
\end{proof}

\begin{theorem}[\lean{Basic.kron\_symm}]
For any $i, j \in \{0,\ldots,n-1\}$, $\delta_{ij} = \delta_{ji}$.
\end{theorem}
\begin{proof}
The Kronecker delta depends symmetrically on its arguments: if $i = j$ then $j = i$, and if $i \ne j$ then $j \ne i$.
\end{proof}

\begin{theorem}[\lean{Basic.d\_kron}]
For any scale algebra, any $k, i, j \in \{0,\ldots,n-1\}$, the Kronecker delta is annihilated by every derivation:
$$\partial_k(\delta_{ij}) = 0.$$
\end{theorem}
\begin{proof}
Since $\delta_{ij} \in \{0, 1\} \subset A$, it is a natural number cast into $A$. By \lean{Basic.ScaleAlgebra.d\_natCast}, every derivation annihilates it.
\end{proof}

\begin{theorem}[\lean{Basic.d\_kron\_mul}]
For any scale algebra, any $k, i, j \in \{0,\ldots,n-1\}$, and any $a \in A$,
$$\partial_k(\delta_{ij} \cdot a) = \delta_{ij} \cdot \partial_k(a).$$
\end{theorem}
\begin{proof}
By the Leibniz rule,
$$\partial_k(\delta_{ij} \cdot a) = \partial_k(\delta_{ij}) \cdot a + \delta_{ij} \cdot \partial_k(a) = 0 + \delta_{ij} \cdot \partial_k(a)$$
since $\partial_k(\delta_{ij}) = 0$ by the previous theorem.
\end{proof}

The following four theorems are the only facts about the Kronecker delta needed by the curvature computation.

\begin{theorem}[\lean{Basic.sum\_kron\_left}]
For any scale algebra, any $b \in \{0,\ldots,n-1\}$, and any function $f : \{0,\ldots,n-1\} \to A$,
$$\sum_{a=0}^{n-1} \delta_{ab} \cdot f(a) = f(b).$$
\end{theorem}
\begin{proof}
The sum contains exactly one nonzero term, when $a = b$. In that case, $\delta_{bb} \cdot f(b) = 1 \cdot f(b) = f(b)$.
\end{proof}

\begin{theorem}[\lean{Basic.sum\_kron\_right}]
For any scale algebra, any $b \in \{0,\ldots,n-1\}$, and any function $f : \{0,\ldots,n-1\} \to A$,
$$\sum_{a=0}^{n-1} \delta_{ba} \cdot f(a) = f(b).$$
\end{theorem}
\begin{proof}
By symmetry of the Kronecker delta, $\delta_{ba} = \delta_{ab}$. The sum then contains exactly one nonzero term when $a = b$, giving $\delta_{bb} \cdot f(b) = f(b)$.
\end{proof}

\subsection{Axioms}

This section develops axioms A2, A3, and A5 of the theory, establishing the scale field and its immediate consequences. The scale $\sigma$ is a dimensionless log-scale with multiplicative representative $s = e^\sigma$ (postulated as a unit of the ring); the scale-energy duality $\varepsilon \cdot s = 1$ is derived from this; and the fiducial covariance (invariance under global shift $\sigma \mapsto \sigma + c$) implies that all geometric objects are exact dimensionless.

\begin{definition}[\lean{Axioms.ScaleField}]
A \emph{scale field} of dimension $n$ on a scale algebra over commutative ring $A$ is a triple $(\sigma, s, \mathbf{d}_s)$ where:
\begin{enumerate}
\item $\sigma \in A$ is the log-scale;
\item $s \in A^\times$ is the scale itself, a unit in the ring;
\item $\mathbf{d}_s$ is the condition that $s$ satisfies the differential equation of the exponential: for each $i \in \{0,\ldots,n-1\}$,
  $$\partial_i(s) = s \cdot \partial_i(\sigma).$$
\end{enumerate}
\end{definition}

\begin{definition}[\lean{Axioms.ScaleField.en}]
Given a scale field $F = (\sigma, s, \mathbf{d}_s)$, the energy scale is defined as the reciprocal of the length scale:
$$\varepsilon := s^{-1}.$$
\end{definition}

\begin{theorem}[\lean{Axioms.ScaleField.en\_mul\_scale}]
For any scale field $F$, the scale-energy duality holds: $\varepsilon \cdot s = 1$.
\end{theorem}
\begin{proof}
Immediate from the definition of $\varepsilon$ as $s^{-1}$ in a commutative ring with units.
\end{proof}

\begin{theorem}[\lean{Axioms.ScaleField.scale\_mul\_en}]
For any scale field $F$, we have also $s \cdot \varepsilon = 1$.
\end{theorem}
\begin{proof}
This follows from $\varepsilon \cdot s = 1$ by commutativity of multiplication in $A$.
\end{proof}

\begin{theorem}[\lean{Axioms.ScaleField.d\_en}]
For any scale field $F$ and any $i \in \{0,\ldots,n-1\}$, the energy scale obeys the dual differential equation:
$$\partial_i(\varepsilon) = -\varepsilon \cdot \partial_i(\sigma).$$
\end{theorem}
\begin{proof}
Differentiate the duality relation $\varepsilon \cdot s = 1$ using the Leibniz rule:
$$0 = \partial_i(\varepsilon \cdot s) = \partial_i(\varepsilon) \cdot s + \varepsilon \cdot \partial_i(s) = \partial_i(\varepsilon) \cdot s + \varepsilon \cdot s \cdot \partial_i(\sigma),$$
using the condition that $\partial_i(s) = s \cdot \partial_i(\sigma)$. Thus
$$\partial_i(\varepsilon) \cdot s = -\varepsilon \cdot s \cdot \partial_i(\sigma).$$
Multiply both sides on the right by $\varepsilon$ (which equals $s^{-1}$):
$$\partial_i(\varepsilon) = -\varepsilon \cdot \partial_i(\sigma).$$
\end{proof}

A global additive shift in the fiducial unit leaves every geometric object constructed from the scale invariant. Let $c \in A$ be a constant (meaning $\partial_i(c) = 0$ for all $i$).

\begin{theorem}[\lean{Axioms.sig\_add\_const}]
The first-order operator on the shifted scale obeys $\partial_i(\sigma + c) = \partial_i(\sigma)$.
\end{theorem}
\begin{proof}
By the Leibniz rule and the fact that $c$ is constant, $\partial_i(\sigma + c) = \partial_i(\sigma) + \partial_i(c) = \partial_i(\sigma) + 0 = \partial_i(\sigma)$.
\end{proof}

\begin{theorem}[\lean{Axioms.hess\_add\_const}]
The second-order operator on the shifted scale obeys $\partial_j(\partial_i(\sigma + c)) = \partial_j(\partial_i(\sigma))$.
\end{theorem}
\begin{proof}
This follows from the previous theorem by applying $\partial_j$ to both sides.
\end{proof}

\begin{theorem}[\lean{Axioms.lap\_add\_const}]
The Laplacian of the shifted scale equals the Laplacian of the original scale: $\Delta(\sigma + c) = \Delta(\sigma)$.
\end{theorem}
\begin{proof}
The Laplacian is $\Delta = \sum_{i=0}^{n-1} \partial_i^2$. Since each second derivative is unchanged by shifting $\sigma$ by a constant, the sum is unchanged.
\end{proof}

\begin{theorem}[\lean{Axioms.gradsq\_add\_const}]
The squared gradient of the shifted scale equals the squared gradient of the original scale: $\sum_{i=0}^{n-1} (\partial_i(\sigma + c))^2 = \sum_{i=0}^{n-1} (\partial_i(\sigma))^2$.
\end{theorem}
\begin{proof}
Since $\partial_i(\sigma + c) = \partial_i(\sigma)$ for each $i$, their squares are equal, and the sum is therefore equal.
\end{proof}

\begin{theorem}[\lean{Axioms.Chr\_add\_const}]
The Christoffel symbols of the shifted scale equal the Christoffel symbols of the original scale: $\Gamma^\rho_{\mu\nu}(\sigma + c) = \Gamma^\rho_{\mu\nu}(\sigma)$.
\end{theorem}
\begin{proof}
The Christoffel symbols depend only on first derivatives of $\sigma$ through the first-order operator. Since $\partial_i(\sigma + c) = \partial_i(\sigma)$, the Christoffel symbols are unchanged.
\end{proof}

\begin{theorem}[\lean{Axioms.Rm\_add\_const}]
The Riemann curvature tensor of the shifted scale equals that of the original scale: $R_{\rho\mu\nu\lambda}(\sigma + c) = R_{\rho\mu\nu\lambda}(\sigma)$.
\end{theorem}
\begin{proof}
The Riemann tensor is built from Christoffel symbols and their derivatives. Since the Christoffel symbols are unchanged, and the derivatives of $\sigma$ are unchanged, the Riemann tensor is unchanged.
\end{proof}

\begin{theorem}[\lean{Axioms.Ric\_add\_const}]
The Ricci curvature tensor of the shifted scale equals that of the original scale: $R_{\mu\nu}(\sigma + c) = R_{\mu\nu}(\sigma)$.
\end{theorem}
\begin{proof}
The Ricci tensor is the trace of the Riemann tensor. Since the Riemann tensor is unchanged, its trace is unchanged.
\end{proof}

\begin{theorem}[\lean{Axioms.Defm\_add\_const}]
The deformation tensor of the shifted scale equals that of the original scale: $D_n(\sigma + c) = D_n(\sigma)$.
\end{theorem}
\begin{proof}
The deformation tensor depends on derivatives of $\sigma$ through the Hessian and first-order operator. All these are unchanged when $\sigma$ is shifted by a constant, so the deformation tensor is unchanged.
\end{proof}

\begin{theorem}[\lean{Axioms.geometry\_fiducial\_invariant}]
\emph{(Axiom A5, collected.)} A global change of the fiducial unit leaves every geometric object invariant. Specifically:
\begin{enumerate}
\item For all $i$, the first-order operator: $\partial_i(\sigma + c) = \partial_i(\sigma)$;
\item For all $\rho, \mu, \nu, \lambda$, the Riemann tensor: $R_{\rho\mu\nu\lambda}(\sigma + c) = R_{\rho\mu\nu\lambda}(\sigma)$;
\item For all $\mu, \nu$, the Ricci tensor: $R_{\mu\nu}(\sigma + c) = R_{\mu\nu}(\sigma)$;
\item For all $i, j$, the deformation tensor: $D_n(\sigma + c)_{ij} = D_n(\sigma)_{ij}$;
\item The bare scalar curvature: $\mathcal{R}(\sigma + c) = \mathcal{R}(\sigma)$.
\end{enumerate}
This encodes that the geometry is exactly dimensionless: no observable depends on where the zero of $\sigma$ is placed. Only scale differences are physical.
\end{theorem}
\begin{proof}
Each statement follows from the corresponding individual theorem above. The bare scalar curvature is defined as the trace of the Ricci tensor, so it inherits invariance from the Ricci tensor.
\end{proof}

The first-order, second-order, and squared-gradient operators are linear in $\sigma$.

\begin{theorem}[\lean{Axioms.sig\_neg}]
The first-order operator satisfies $\partial_i(-\sigma) = -\partial_i(\sigma)$.
\end{theorem}
\begin{proof}
By the Leibniz rule applied to negation (which is the scalar multiplication by $-1$), and the fact that constants are annihilated by derivations.
\end{proof}

\begin{theorem}[\lean{Axioms.hess\_neg}]
The second-order operator satisfies $\partial_j(\partial_i(-\sigma)) = -\partial_j(\partial_i(\sigma))$.
\end{theorem}
\begin{proof}
Applying $\partial_j$ to $\partial_i(-\sigma) = -\partial_i(\sigma)$ preserves the sign.
\end{proof}

\begin{theorem}[\lean{Axioms.lap\_neg}]
The Laplacian satisfies $\Delta(-\sigma) = -\Delta(\sigma)$.
\end{theorem}
\begin{proof}
The Laplacian is a sum of second derivatives, each of which satisfies $\partial_i^2(-\sigma) = -\partial_i^2(\sigma)$, so the sum reverses sign.
\end{proof}

\begin{theorem}[\lean{Axioms.gradsq\_neg}]
The squared gradient satisfies $\sum_{i=0}^{n-1} (\partial_i(-\sigma))^2 = \sum_{i=0}^{n-1} (\partial_i(\sigma))^2$.
\end{theorem}
\begin{proof}
Since $(\partial_i(-\sigma))^2 = (-\partial_i(\sigma))^2 = (\partial_i(\sigma))^2$, the squared gradient is unchanged when $\sigma$ is negated.
\end{proof}

\subsection{Postulates}

This section collects the seven postulates of Scale-Coupled Dynamics in one place, with their formal carriers and status. Of the seven, three (A2, A3, A5) are developed in Axioms; four (A1, A4, A6, A7) are carriers elsewhere. Here we verify that Axiom A1 has a single formal carrier (not two), restate Axiom A7 where the others appear, and derive key consequences.

Axiom A1 is the substrate: a bare counting structure carrying $n$ derivations that commute as operators. This is declared both as a commutative-ring specialization (\textsc{ScaleAlgebra}) and as a ring specialization (\textsc{DiffRing}). The following theorem and definition prove they are the same.

\begin{definition}[\lean{Postulates.scaleAlgebra\_is\_diffRing}]
Every scale algebra (commutative ring with $n$ commuting derivations) gives rise to a DiffRing: the same derivations, the same Leibniz rule, the same commutativity of mixed partials. The construction is formal identity.
\end{definition}

\begin{theorem}[\lean{Postulates.scaleAlgebra\_is\_diffRing\_d}]
Under the above identification, the derivations are preserved: for any $i \in \{0,\ldots,n-1\}$ and $a \in A$, the $i$-th derivation in the DiffRing is exactly $\partial_i(a)$ in the scale algebra.
\end{theorem}
\begin{proof}
This is immediate from the definition, as the two structures carry the same derivations.
\end{proof}

Axiom A7 (Observability) states that the number of rotations generated by the coupling equals the number of scale directions. Formally, this is the condition that $2 \cdot \dim(\Lambda^2 \mathbb{R}^k) = 2 \cdot k$ in doubled form to avoid division, where $k = m + 1$ and $m$ parameterizes the dimension of the defining representation.

\begin{definition}[\lean{Postulates.Observability}]
The observability condition at level $m$ is that $(m+1) \cdot m = 2 \cdot (m+1)$.
\end{definition}

\begin{theorem}[\lean{Postulates.three\_from\_A7}]
If the observability condition holds at level $m$, then $m = 2$.
\end{theorem}
\begin{proof}
From $(m+1) \cdot m = 2 \cdot (m+1)$, we can rewrite using commutativity as $m \cdot (m+1) = 2 \cdot (m+1)$. Since $m + 1 \ge 1$ is nonzero, we may cancel to obtain $m = 2$.
\end{proof}

\begin{theorem}[\lean{Postulates.A7\_holds\_only\_at\_three}]
The observability condition holds if and only if $m = 2$.
\end{theorem}
\begin{proof}
The forward direction is the previous theorem. For the reverse, if $m = 2$, then $(m+1) \cdot m = 3 \cdot 2 = 6 = 2 \cdot 3 = 2 \cdot (m+1)$ by direct calculation.
\end{proof}
